% ============================================================================
% Preâmbulo: classe, pacotes e dados do trabalho
% ============================================================================

\PassOptionsToPackage{hyperfootnotes=false}{hyperref}
\documentclass[
    12pt,
    openright,
    oneside,
    a4paper,
    sumario=tradicional,
    brazilian,
    english,
]{abntex2}

% Idioma do texto: português (padrão). Para escrever em inglês, descomente a linha abaixo.
% \def\textoemingles{}

\usepackage{abntex2-UFV}

% Fonte, codificação e espaçamento
\usepackage[T1]{fontenc}
\usepackage[utf8]{inputenc}
\usepackage{newtxtext,newtxmath}   % Times
\usepackage{indentfirst}
\OnehalfSpacing

% Figuras, tabelas e matemática
\usepackage{graphicx}
\usepackage{float}
\usepackage{booktabs}
\usepackage{multirow}
\usepackage{subcaption}
\usepackage{caption}
\usepackage{amsmath}
\usepackage{mathtools}
\usepackage{bm}
\usepackage{tikz}
\usepackage{xurl}

% Ambiente de definição numerado por capítulo
\let\openbox\relax
\usepackage{amsthm}
\theoremstyle{definition}
\newtheorem{definicao}{Definição}[chapter]

% Citações numéricas no padrão ABNT, com o número entre colchetes
\usepackage[num,abnt-emphasize=bf]{abntex2cite}
\citebrackets{[}{]}

% Texto em inglês: títulos das listas, do sumário e das legendas em inglês
\ifdefined\textoemingles
    \addto\captionsbrazil{%
        \def\bibname{References}%
        \renewcommand{\figurename}{Figure}%
        \renewcommand{\tablename}{Table}%
        \renewcommand{\listfigurename}{List of Figures}%
        \renewcommand{\listtablename}{List of Tables}%
        \renewcommand{\contentsname}{Contents}%
        \renewcommand{\chaptername}{Chapter}%
    }
    \addto\captionsenglish{\def\bibname{References}}
    \AtBeginDocument{\selectlanguage{english}}
\fi

% Títulos de capítulos e seções sem fonte sans-serif
\renewcommand{\ABNTEXchapterfont}{\rmfamily\bfseries}
\renewcommand{\ABNTEXsectionfont}{\rmfamily\bfseries}
\renewcommand{\ABNTEXsubsectionfont}{\rmfamily\bfseries}
\renewcommand{\ABNTEXsubsubsectionfont}{\rmfamily\bfseries}

% ----------------------------------------------------------------------------
% Dados do trabalho
% ----------------------------------------------------------------------------
\titulo{\normalsize{\textbf{Título da Dissertação: Subtítulo}}}
\autor{\normalsize{\textbf{Nome Completo do Autor}}}
\local{\normalsize{\textbf{Florestal - Minas Gerais}}}
\data{\normalsize{\textbf{2026}}}
\orientador{Nome do Orientador}
% \coorientador{Nome do Coorientador}
\instituicao{Universidade Federal de Viçosa}
\campus{Campus Florestal}
\curso{Pós-graduação em Ciência da Computação}
\preambulo{Dissertação apresentada à Universidade Federal de Viçosa, como parte das
exigências do Programa de Pós-Graduação em Ciência da Computação, para obtenção do título
de \textit{Magister Scientiae}.}

% Metadados e cores dos links no PDF
\makeatletter
\hypersetup{
    pdftitle={\@title},
    pdfauthor={\@author},
    colorlinks=true,
    linkcolor=black,
    citecolor=blue,
    urlcolor=blue,
    bookmarksdepth=4,
}
\makeatother
